\chapter{把诊断、尺度和误差算到底}
\label{tm:chap:worked}

本章每个算例都给出输入、代入和判断。前三例采用自设教学参数，后面几例明确指出哪些输入摘自实验论文、哪些是为说明方法另加的近似。它们不是对原始实验数据的重新拟合。

\section{算例一：法向磁扰动怎样换成磁岛全宽}

取局地平衡和单一扰动
\begin{equation}
 \psi(x,y)=\frac{B'x^2}{2}+\widetilde\psi\cos ky,
 \qquad B'=2\,\mathrm{T\,m^{-1}},\quad k=20\,\mathrm{m^{-1}}.
\end{equation}
按本讲义约定 $\B=\widehat{\vct z}\times\nabla\psi$，法向分量为
\begin{equation}
 \delta B_x=-\partial_y\psi=k\widetilde\psi\sin ky.
\end{equation}
假设在有理面附近测得该单模法向场峰值 $b_x=0.20\,\mathrm{mT}$，则
\begin{equation}
 \widetilde\psi=\frac{b_x}{k}
 =\frac{2.0\times10^{-4}\,\mathrm T}{20\,\mathrm{m^{-1}}}
 =1.0\times10^{-5}\,\mathrm{T\,m}.
\end{equation}
量纲检查：二维磁通函数是 $\mathrm{T\,m}$，不是总磁通 $\mathrm{Wb}$，因为我们没有再乘出平面的长度。

X 点在 $(x,ky)=(0,0)$，分离界值为 $\psi_X=\widetilde\psi$。通过 O 点的径向截线是 $ky=\pi$。分离界与此截线相交时，
\begin{align}
 \frac{B'x^2}{2}-\widetilde\psi&=\widetilde\psi,\\
 x_{\rm sep}^2&=\frac{4\widetilde\psi}{B'},\\
 W=2|x_{\rm sep}|&=4\sqrt{\frac{\widetilde\psi}{B'}}
 =4\sqrt{\frac{10^{-5}}{2}}\,\mathrm m
 =8.94\,\mathrm{mm}.
 \label{tm:worked_width}
\end{align}
因此半宽为 $4.47\,\mathrm{mm}$。若误把 $W$ 写成半宽，后续涉及 $W^2$ 的项会差四倍。还要注意：这是\emph{有理面附近}的扰动幅度。壁上线圈的 $0.20\,\mathrm{mT}$ 不应直接代入；必须用外区本征函数、几何和壁响应把信号映射到共振面。

\section{算例二：有理面位置与平行相位尺度}

使用图\ref{tm:fig:rational}的自设剖面
\begin{equation}
 q(r)=1+2(r/a)^2,\qquad a=0.60\,\mathrm m,\qquad R=1.70\,\mathrm m.
\end{equation}
求 $2/1$ 面：$1+2(r_s/a)^2=2$，所以
\begin{equation}
 r_s=\frac{a}{\sqrt2}=0.4243\,\mathrm m.
\end{equation}
对 $3/2$ 面，$1+2(r_s/a)^2=1.5$，得到 $r_s=a/2=0.3000\,\mathrm m$。注意 $q=2$ 是磁力线绕行比，不是“半径等于 2”。

对于 $m=2,n=1$，在环向场占优的近似 $B_\zeta/B\simeq1$ 下，
\begin{align}
 q'(r_s)&=\frac{4r_s}{a^2}=4.714\,\mathrm{m^{-1}},\\
 k_\parallel'(r_s)&=-\frac{m q'}{R q_s^2}
 =-\frac{2\times4.714}{1.70\times4}
 =-1.386\,\mathrm{m^{-2}}.
\end{align}
离共振面向外 $x=1.0\,\mathrm{cm}$ 时，一阶估计给出
\begin{equation}
 k_\parallel\simeq-1.386\times0.010
 =-1.386\times10^{-2}\,\mathrm{m^{-1}},
 \qquad L_{\parallel,\alpha}\equiv\frac{2\pi}{|k_\parallel|}\simeq453\,\mathrm m.
\end{equation}
这是沿磁力线使\emph{扰动相位}改变 $2\pi$ 所需的局地长度，不是托卡马克的周长，也不是任何热粒子实际自由飞行的平均自由程。共振面上该相位长度发散，说明相位与平衡场线相锁；磁力线仍然在强磁场中延伸。

这一结果还解释了为什么小岛未必立即完全均温。仅作教学估计，取 $\chi_\parallel=10^8\,\mathrm{m^2/s}$、$\chi_\perp=1\,\mathrm{m^2/s}$，并\emph{暂用} $453\,\mathrm m$ 作为平行长度，则
\begin{equation}
 \tau_\parallel\sim\frac{453^2}{10^8}=2.05\,\mathrm{ms},
 \qquad
 \tau_\perp\sim\frac{(0.010)^2}{1}=0.10\,\mathrm{ms}.
\end{equation}
即使 $\chi_\parallel/\chi_\perp$ 很大，长连接距离也可能抵消它。这不是严格求出的磁岛平坦化阈值：实际岛内应沿岛磁场求连接长度和二维热输运。它说明比较扩散\emph{系数}之前，还必须比较扩散\emph{距离的平方}。

\section{算例三：从相位、幅度和电压得到不同物理量}

\subsection{先定模数，再定频率}

设一组理想合成信号满足式\eqref{tm:mirnov_B}。在相同环向角上，两个探头相差 $\Delta\theta=\pi/4$，展开后的相位差为 $\pi/2$。因此
\begin{equation}
 m=\frac{\Delta\alpha}{\Delta\theta}=2.
\end{equation}
另一对探头有相同极向角，$\Delta\zeta=\pi/3$，相位差为 $-\pi/3$，得到
\begin{equation}
 n=-\frac{\Delta\alpha}{\Delta\zeta}=1.
\end{equation}
如果只保留模 $2\pi$ 的相位而没有更多探头，第一对还可与其他整数相混淆。例如 $m=10$ 也给出 $10\pi/4=\pi/2+2\pi$。完整阵列及可信的模数范围才消除这种混叠。

周期测得 $T=0.10\,\mathrm{ms}$，故 $f=1/T=10\,\mathrm{kHz}$，$\omega=2\pi f=6.283\times10^4\,\mathrm{s^{-1}}$。频率表示模式旋转，不是增长率。旋转方向还需按照探头坐标和相位约定确定，单一功率谱峰不包含符号。

\subsection{磁场幅度增长四倍，岛宽只增长两倍}

假设已校正为磁场的单模包络在 $\Delta t=0.20\,\mathrm{ms}$ 内由 $0.10$ 增到 $0.40\,\mathrm{mT}$，背景及 $k$ 保持不变，并且这段确实处在线性阶段。由于 $b_x=k\widetilde\psi$，
\begin{equation}
 \gamma=\frac{\ln[b_x(t_2)/b_x(t_1)]}{t_2-t_1}
 =\frac{\ln4}{2.0\times10^{-4}}
 =6.931\times10^3\,\mathrm{s^{-1}},
 \quad \gamma^{-1}=144.3\,\mu\mathrm s.
\end{equation}
因 $W\propto\sqrt{\widetilde\psi}$，$W_2/W_1=\sqrt4=2$，岛宽对数斜率为
\begin{equation}
 \frac{\dd\ln W}{\dd t}=\frac\gamma2
 =3.466\times10^3\,\mathrm{s^{-1}}.
\end{equation}
把岛宽半对数图的斜率直接当 $\gamma$，会少一倍。

如果两个磁场幅度各有独立的 $5\%$ 标准不确定度，暂忽略时钟误差，对数误差传播给出
\begin{equation}
 \sigma_\gamma\simeq\frac{1}{\Delta t}
 \sqrt{\left(\frac{\sigma_{b_1}}{b_1}\right)^2+
       \left(\frac{\sigma_{b_2}}{b_2}\right)^2}
 =\frac{\sqrt{0.05^2+0.05^2}}{2\times10^{-4}}
 =354\,\mathrm{s^{-1}}.
\end{equation}
结果为 $(6.93\pm0.35)\times10^3\,\mathrm{s^{-1}}$，约 $5.1\%$。共同的乘法标定误差会在幅度比中抵消，不能总当作两次独立误差相加；平衡演化和噪声底则是另外的系统问题。

\subsection{为什么不能直接用线圈电压比}

若旋转频率从 $10$ 降到 $5\,\mathrm{kHz}$，真实磁场幅度同时增长四倍，则在 $|\omega|\gg\gamma$ 的近似下，电压只增长约
\begin{equation}
 \frac{|V_2|}{|V_1|}\simeq\frac{f_2b_2}{f_1b_1}=2.
\end{equation}
直接代入电压比会给出 $\ln2/\Delta t=3.466\times10^3\,\mathrm{s^{-1}}$，约少一半。若保留本例有限 $\gamma$，应使用 $\sqrt{\omega^2+\gamma^2}$，电压比约为 $2.04$，误差仍然很大。这一例不是某次放电的再分析，而是说明仪器测量方程本身就参与物理解释。

\section{算例四：MRX 为什么可以既有 Hall 效应又有电阻}

取 \citet{TMJaraAlmonte2016} 的代表性测量尺度，见式\eqref{tm:mrx_scales}。直接计算
\begin{align}
 \frac{\delta}{d_i}&=\frac9{200}=0.045,\qquad
 \frac{\delta}{d_e}=\frac9{0.5}=18,\\
 \frac{\lambda_{\rm mfp}}{\delta}&=\frac29=0.222,\qquad
 \frac{\Delta x_{\rm probe}}{\delta}=\frac69=0.667.
\end{align}
第一比值告诉我们片远薄于离子惯性长度，忽略离子电子相对漂移不可信。第二比值说明片整体比电子惯性长度大。第三比值说明跨片距离可包含数个电子平均自由程，碰撞电阻并不自动消失。它不是强渐近极限，仍要检查局地欧姆定律；原文还以电场/电流的有效电阻与 Spitzer 值作比较。第四比值提醒我们片半厚度仅相当于约 $1.5$ 个最好跨片采样间距，不能声称这组探针解析了所有电子细结构。

原文给出的典型 Alfv\'en 时间为 $8.8\pm3\,\mu\mathrm s$，每次用于统计的一段观测时间约 $6.4\,\mu\mathrm s$。用中心值计算
\begin{equation}
 \frac{6.4}{8.8}=0.73.
\end{equation}
这是统计窗口与系统特征时间的比较，不是仪器带宽，也不是说每个岛只有一个时间采样点。研究用许多放电形成统计集合；事件平均的误差可以比事件之间的实际散布小。若目的是测一个缓慢线性本征模的 $\gamma$，还必须展示单个事件中足够长且平衡稳定的指数增长区间。

这些尺度不能擅自以“氢等离子体”公式重算：该实验使用氩，$d_i$ 依赖离子质量、电荷态和密度。这里使用论文报告的 $d_i$，没有猜测氩的平均电荷态来人为提高参数精度。

\section{算例五：激光实验里，“未看见增长”是否等于“不增长”}

\subsection{先统一单位，再比较电阻尺度}

从 \citet{TMPearcy2024} 表 I 取代表值
\begin{equation}
 L=2000\,\mu\mathrm m=2.0\times10^{-3}\,\mathrm m,
 \qquad V_A=1.5\times10^{11}\,\mu\mathrm{m/s}
 =1.5\times10^5\,\mathrm{m/s},
\end{equation}
\begin{equation}
 \eta\simeq4.0\times10^{-2}\,\Omega\,\mu\mathrm m
 =4.0\times10^{-8}\,\Omega\,\mathrm m.
\end{equation}
不要把 $\Omega\,\mu\mathrm m$ 误读为 $\Omega\,\mathrm m$，否则差一百万倍。磁扩散率为
\begin{equation}
 D_\eta=\frac{\eta}{\mu_0}
 =\frac{4.0\times10^{-8}}{4\pi\times10^{-7}}
 =3.183\times10^{-2}\,\mathrm{m^2/s}.
\end{equation}
若直接以所列 $L$ 作特征长度，
\begin{align}
 S_L&=\frac{LV_A}{D_\eta}
 =\frac{(2.0\times10^{-3})(1.5\times10^5)}{3.183\times10^{-2}}
 =9.42\times10^3,\\
 \tau_{AL}&=\frac L{V_A}=1.33\times10^{-8}\,\mathrm s=13.3\,\mathrm{ns}.
\end{align}
这一 $S$ 的数量级与论文范围一致。若别人把同一几何的半长定义为 $L$，必须先转换，不能争论一个相差二的数字。

\subsection{电流全半高宽不是 Harris 半厚度}

为了演示尺度比较，形式上代入稳态 Sweet--Parker 估计
\begin{equation}
 a_{\rm SP}\sim\frac L{\sqrt{S_L}}
 =20.6\,\mu\mathrm m.
\end{equation}
这只是参照尺度：该实验快速形成、强驱动且接近离子尺度，不能据此认定它实际处于稳态 Sweet--Parker 平衡。

Harris 平衡满足 $B_y=B_0\tanh(x/a)$，于是
\begin{equation}
 J_z=\frac{B_0}{\mu_0a}\operatorname{sech}^2(x/a).
\end{equation}
半峰值处 $\operatorname{sech}^2(x_{1/2}/a)=1/2$，所以
\begin{equation}
 \cosh(x_{1/2}/a)=\sqrt2,\qquad
 w_{J,\rm FWHM}=2x_{1/2}=2a\operatorname{arcosh}\sqrt2=1.76275a.
\end{equation}
假定为演示而取 $w_{J,\rm FWHM}=45\,\mu\mathrm m$，则 $a=25.5\,\mu\mathrm m$，$a/d_i=25.5/55=0.464$。这里“45”是论文所报宽度范围内的\emph{教学取值}，不是某时刻数据的重新提取。实际轮廓未必严格 Harris，因此换算也是模型条件下的换算。

\subsection{把增长速度乘上真实观察时间}

论文讨论约 $0.885$ 至 $0.952\,\mathrm{ns}$ 的若干照相时刻，总窗约 $\Delta t=0.067\,\mathrm{ns}$，空间分辨率约 $5\,\mu\mathrm m$。设一个\emph{供比较的}局地速度为 $v_*=0.1V_A$，则
\begin{equation}
 v_*=1.5\times10^4\,\mathrm{m/s}=15\,\mu\mathrm{m/ns},
 \qquad v_*\Delta t=15\times0.067=1.01\,\mu\mathrm m.
\end{equation}
即使宽度的两侧都以此速率外移，全宽增加也只有约 $2.0\,\mu\mathrm m$，仍小于 $5\,\mu\mathrm m$。单侧位移达到该分辨率需要
\begin{equation}
 \Delta t_{\rm resolve}\sim\frac{5}{15}\,\mathrm{ns}=0.333\,\mathrm{ns},
\end{equation}
约为现有总窗的五倍。此处 $v_*$ 不是由 $\gamma W/2$ 测出来的岛边速度，而是原文讨论所用典型重联速度数量级的敏感性比较。它表明图 6 中误差内没有增长，与物理上存在缓慢演化完全相容。

相反，若某个边缘结构表观扩展到约 $1300\,\mu\mathrm{m/ns}$，这与膨胀泡 $500$--$800\,\mu\mathrm{m/ns}$ 的流动和几何变化相当，不能不加区分地除以岛宽得到线性增长率。\emph{没有分辨出变化}和\emph{变化很快但不是所需本征过程}，是两种不同的实验限制。\citep{TMPearcy2024}

\section{算例六：ECCD 偏移一厘米为何可能很重要}

\citet{TMPetty2003} 的射线追踪给出代表性电流沉积全半高宽 $w_{\rm CD}=2.6\,\mathrm{cm}$，反馈调整对应位置移动约 $0.9\,\mathrm{cm}$。为单独展示几何敏感性，设一个\emph{自选高斯轮廓}
\begin{equation}
 j_{\rm CD}(r)=j_0\exp\left[-4\ln2\frac{(r-r_{\rm CD})^2}{w_{\rm CD}^2}\right].
\end{equation}
这个系数保证 $r-r_{\rm CD}=\pm w_{\rm CD}/2$ 时恰好半峰值。若沉积中心与有理面差 $\Delta r=0.9\,\mathrm{cm}$，该面上的局地电流相对峰值为
\begin{equation}
 \frac{j_{\rm CD}(r_s)}{j_0}
 =\exp\left[-4\ln2\left(\frac{0.9}{2.6}\right)^2\right]
 =0.717.
\end{equation}
仅这个偏移就降低约 $28\%$。若要求该局地值至少是峰值的 $80\%$，解不等式
\begin{equation}
 |\Delta r|\leq w_{\rm CD}\sqrt{\frac{-\ln0.8}{4\ln2}}
 =0.738\,\mathrm{cm}.
\end{equation}
因此一厘米在大装置中看起来很短，相对于厘米级沉积宽度却已经重要。

再看展宽。在总驱动电流和局地几何近似不变时，高斯面积 $\int j_{\rm CD}\dd r\propto j_0w_{\rm CD}$，所以 $j_0\propto1/w_{\rm CD}$。宽度从 $2.6$ 变到 $3.3\,\mathrm{cm}$ 时，峰值变为原来的
\begin{equation}
 \frac{j_{0,2}}{j_{0,1}}=\frac{2.6}{3.3}=0.788.
\end{equation}
原文无展宽的峰值 $j_{\rm CD}/j_{\rm BS}\simeq2.5$，按此固定总电流近似会降到约 $1.97$。真实稳定作用还要将驱动与岛的螺旋结构重叠积分，并考虑旋转与调制。因此 $0.717$ 不是实际抑制效率，$0.738\,\mathrm{cm}$ 也不是 DIII-D 经实验验证的普适容差；它们是清楚标明假设的局地几何算例。

在这篇报告采用的完美对齐、非调制模型中，还有一个随宽度变化的效率因子
\begin{equation}
 \eta_{\rm CD}=\frac{0.4}{1+2w_{\rm CD}^2/W^2}.
\end{equation}
这里 $\eta_{\rm CD}$ \textbf{无量纲}，绝不是本讲义的电阻率 $\eta$。对部分抑制的 $W=8\,\mathrm{cm}$，$w_{\rm CD}=2.6$ 和 $3.3\,\mathrm{cm}$ 分别给出 $0.330$ 与 $0.298$。峰值降低与重叠效率变化是两个相关影响，不能只算一个比值就宣称复现了完整修正 Rutherford 计算。\citep{TMPetty2003}

\section{最后把证据链写成一句完整的话}

面对一组新结果，一个有用的表达形式是：“在给定平衡和时间区间内，诊断直接测得某个信号，经明确的转换得到某个扰动量；它在误差范围内支持某个模型的某一项预测。”例如本章合成 Mirnov 算例得到的是磁通幅度的增长率，不是未经处理的电压增长率；MRX 比较支持碰撞性双流体尺度排序，不是宣称电子扩散层已经全部解析；激光算例得到的是分辨率限制，不是零增长；DIII-D 算例说明对齐和展宽的物理敏感性，不是替代完整装置控制模型。这些区分正是把推导真正用于实验时最需要掌握的部分。
